% =============================================================================
% CAPÍTULO 3 — METODOLOGÍA
% =============================================================================
\chapter{Metodología}
\label{ch:metodologia}

% ---------------------------------------------------------------------------
% NOTA: Este capítulo está estructurado con marcadores TODO para completar.
% El contenido de ejemplo anterior fue reemplazado por la estructura definitiva.
% ---------------------------------------------------------------------------

\section{Paradigma y enfoque de investigación}
\label{sec:enfoque}

% TODO: Describir el enfoque de investigación adoptado:
%   - Design Science Research (DSR) como marco metodológico principal
%     (Hevner et al., 2004): el artefacto de software ES la contribución
%   - Justificar por qué DSR es adecuado para este problema
%     (problema práctico + sistema funcional + evaluación empírica)
%   - Mencionar el carácter aplicado y constructivo del trabajo
%   - Posicionamiento cuali-cuantitativo: análisis documental + métricas
%     de desempeño del sistema
\ldots

\section{Preguntas de investigación y objetivos operativos}
\label{sec:preguntas_investigacion}

% TODO: Reformular las preguntas de investigación en términos operacionales:
%   - PG: ¿Puede un sistema MAS+LLM+MCP asistir la gestión documental
%     para acreditación universitaria con trazabilidad y cumplimiento normativo?
%   - PE1: ¿Qué información clave requiere la CONEAU y cómo se modela?
%   - PE2: ¿Cómo se estructura la arquitectura multi-agente para el dominio?
%   - PE3: ¿Qué nivel de precisión alcanza la extracción automática de propuestas?
%   - PE4: ¿En qué medida MCP mejora la interoperabilidad y trazabilidad
%     respecto a una arquitectura sin protocolo estandarizado?
\ldots

\section{Proceso de desarrollo adoptado}
\label{sec:proceso_desarrollo}

% TODO: Describir el ciclo de desarrollo iterativo-incremental:
%   - Etapa 1 — Relevamiento: análisis de propuestas docentes reales de la
%     carrera \programaTesis; revisión de Ordenanza 75/23 y estándares CONEAU;
%     entrevistas/consultas con docentes y secretaría académica
%   - Etapa 2 — Modelado: definición del modelo de datos, mapa de roles
%     (docente, director, secretaría, par evaluador) y flujo de trabajo
%   - Etapa 3 — Diseño arquitectónico: selección MAS+LLM+MCP+RAG;
%     definición de agentes y sus responsabilidades
%   - Etapa 4 — Implementación: backend FastAPI, frontend React, agentes,
%     módulos de extracción, búsqueda semántica
%   - Etapa 5 — Evaluación: pruebas sobre corpus de propuestas reales;
%     métricas de extracción + evaluación de asistencia LLM
%   - Etapa 6 — Refinamiento: ajuste de agentes, prompts y flujos
%     a partir de resultados
%   Incluir un diagrama de línea de tiempo o proceso (TikZ o figura)
\ldots

\section{Universo, corpus y fuentes de datos}
\label{sec:corpus}

% TODO: Describir detalladamente el corpus de trabajo:
%   - Carrera: \programaTesis — cuántas asignaturas, cuántas propuestas
%   - Origen de los documentos: secretaría académica UCASAL
%   - Formatos: DOCX (formato institucional), Google Docs
%   - Variabilidad del corpus: distintas versiones del template,
%     distintos años lectivos, distintos docentes
%   - Criterios de inclusión/exclusión de documentos
%   - Proceso de anonimización (si aplica)
%   - Cómo se construyó el ground truth para evaluación
\ldots

\section{Métricas de evaluación del sistema}
\label{sec:metricas}

% TODO: Definir las métricas usadas para cada dimensión:
%   - Extracción de información: Precisión, Recall, F1 por campo
%     (título, fundamentación, objetivos, contenidos, RA, bibliografía)
%   - Asistencia LLM: BLEU/ROUGE para generación de texto;
%     evaluación humana (escala Likert) para calidad de sugerencias
%   - Trazabilidad MCP: % de intercambios logueados correctamente;
%     verificación de completitud del log de auditoría
%   - Usabilidad: métricas SUS (System Usability Scale) o cuestionario propio
%     si se realizaron pruebas con usuarios
%   Incluir las fórmulas formales (align)
\ldots

\section{Herramientas, tecnologías y entorno de desarrollo}
\label{sec:herramientas}

% TODO: Completar la tabla de tecnologías usadas con versiones reales:
%   - Backend: FastAPI + SQLAlchemy + SQLite/PostgreSQL
%   - Frontend: React 18 + Vite 5
%   - Extracción DOCX: python-docx
%   - LLM: OpenAI GPT-4o / GPT-4 (indicar modelo exacto usado en producción)
%   - Embeddings: text-embedding-3-small o ada-002 (indicar cuál)
%   - Vector store: hnswlib o ChromaDB (indicar cuál quedó en producción)
%   - Agentes: LangChain / LlamaIndex / agentes propios (indicar)
%   - Protocolo MCP: SDK Python de Anthropic (versión)
%   - Control de versiones: Git + GitHub
%   - Entorno: Python 3.11+, Node.js 20+
%   Armar tabla formal \ref{tab:tecnologias}
\ldots

\section{Consideraciones éticas y de privacidad}
\label{sec:etica}

% TODO: Mencionar brevemente:
%   - Los documentos docentes son de naturaleza institucional,
%     no personal — no aplican restricciones GDPR/PDPA directas
%   - Si se usaron propuestas reales: consentimiento institucional,
%     anonimización de datos de alumnos si los hubiera
%   - Uso de APIs externas (OpenAI): datos enviados a servicios de terceros;
%     política de no-training de OpenAI para API tier
%   - Reproducibilidad: versiones fijas de modelos, semillas fijas
\ldots
